\hypertarget{appendix-2391db8052b1}{}\label{referencestart:appendix-2391db8052b1}
\section{Cross products and oriented area}\label{appendix-2391db8052b1-cross-products-and-oriented-area}

\textbf{Curriculum depth:} Introduced · \textbf{First seen:} Phase6 · \textbf{Latest development:} Phase6

\begin{allmightroledefinition}

\subsection{Definition and mechanism}\label{appendix-2391db8052b1-definition-and-mechanism}

In oriented Euclidean three-dimensional space, the cross product is a perpendicular area vector. Its magnitude is parallelogram area; swapping inputs reverses its orientation. Ordinary component formulas assume an oriented orthonormal basis.

\end{allmightroledefinition}

Panel edges \((4,0,0)\) and \((0,3,0)\) give \((0,0,12)\), so triangle area is \(6\). This is also the reconstructed Part 542 definition missing as a standalone supplied lesson.

\subsection{Prerequisites}\label{appendix-2391db8052b1-prerequisites}

\begin{itemize}
\tightlist
\item
  \href{../06-GPU-Architecture-and-Real-Time-Rendering/main.pdf\#nameddest=appendix-8bc4e8fc547c}{Dot products and orthogonality (GPU Architecture and Real-Time Rendering, Reference notes: Dot products and orthogonality)} --- Inner-product measurements supply the perpendicularity and alignment constraints.
\end{itemize}

\subsection{Used by}\label{appendix-2391db8052b1-used-by}

\begin{itemize}
\item
  \href{../02-Mathematics-for-Visual-Computing/main.pdf\#nameddest=appendix-da478be01d02}{Determinants orientation and volume (Mathematics for Visual Computing, Reference notes: Determinants orientation and volume)}
\item
  \href{../02-Mathematics-for-Visual-Computing/main.pdf\#nameddest=appendix-965d9260eeab}{Quaternion algebra and Hamilton product (Mathematics for Visual Computing, Reference notes: Quaternion algebra and Hamilton product)}
\item
  \hyperref[appendix-87d7626c8d37]{Coordinate handedness and orientation} --- Builds on this mechanism in the connected geometry story.
\item
  \hyperref[appendix-5caa058bd9c7]{Triple products and signed volume} --- Builds on this mechanism in the connected geometry story.
\item
  \hyperref[appendix-e569e60e18e6]{Robust vector angles} --- Builds on this mechanism in the connected geometry story.
\item
  \hyperref[appendix-e31893dec985]{Surface tangents normals and bitangents} --- Builds on this mechanism in the connected geometry story.
\item
  \hyperref[appendix-341296fb3f64]{Robust geometric predicates} --- Builds on this mechanism in the connected geometry story.
\item
  \hyperref[appendix-e1960925d1b8]{Exterior products and geometric algebra} --- Builds on this mechanism in the connected geometry story.
\item
  \href{../03-Geometry-Cameras-and-Image-Formation/main.pdf\#nameddest=appendix-7107009c505b}{Local rigidity and mesh orientation fields (Geometry, Cameras, and Image Formation, Reference notes: Local rigidity and mesh orientation fields)} \#\# Recall and next depth
\end{itemize}

Prove perpendicularity and the area identity; distinguish a two-dimensional signed area from a vector cross. Full production proofs, target-specific integration and learner mastery require further evidence.

\subsection{Source and navigation}\label{appendix-2391db8052b1-source-and-navigation}

\href{../02-Mathematics-for-Visual-Computing/main.pdf\#nameddest=01-foundations-i-ch017-oriented-area-tells-us-which-side-of-the-panel-is-front}{Teaching narrative (Mathematics for Visual Computing, Ch. 1)} · \hyperref[appendix-fd090d982a15-phase-6---parts-501-600]{Coverage and source gap} · \hyperref[appendix-a103e7846d21-phase-6---source-analysis-and-verification]{Evidence and limits} · \hyperref[appendix-4ad47bdd269a]{Dependency Map}
